\section{Introduction} \label{sec:introduction}
%The ability to quickly analyze the overwhelming amounts of sequencing data enables several scientific advancements in personalized medicine and understanding genomic contributions to health across population \citep{hindorff2018prioritizing}. 
One of the most fundamental computational steps in most genomic analyses is {\emph{sequence alignment}}~\citep{alser2020technology, senol2019nanopore}. 
This step is formulated as \alser{an} \emph{approximate string matching} (ASM) problem \citep{navarro2001guided} and \alser{it calculates: (1) \emph{edit distance} between two given sequences, (2) type of each edit (i.e., insertion, deletion, or substitution), and (3) location of each edit in one of the two given sequences.}
Edit distance is defined as the minimum number of edits needed to convert one sequence into the other \citep{levenshtein1966binary}. \alser{These edits result from both sequencing errors \citep{firtina2020apollo} and genetic variations \citep{10002015global}. }
Edits can have different weights, based on a user-defined \emph{scoring} function, to allow favoring one edit type over \alser{another} \citep{wang2011comparison}. 
Sequence alignment involves a \emph{backtracking step}, which calculates an ordered list of characters representing the location and type of each possible edit operation required to change one of the two given sequences into the other. 
As any two sequences can have several different arrangements of the edit operations, we need to examine all possible \emph{prefixes} of the two input sequences and keep track of the pairs of prefixes that provide a minimum edit distance. 
%Enumerating all possible prefixes is necessary for tolerating edits.
Therefore, sequence alignment approaches are typically implemented as dynamic programming (DP) algorithms to avoid re-examining the same prefixes {many times} \citep{alser2020technology,eddy2004dynamic}. 
DP-based \alser{sequence alignment algorithms}, such as Needleman-Wunsch \citep{needleman1970general}, are computationally expensive as they have quadratic time and space complexity (i.e., O($m^2$) for a sequence length of $m$). Many attempts were made to boost the performance of existing sequence aligners. Recent attempts tend to follow one of two \alser{key directions}{, as we comprehensively survey} in~\citep{alser2020accelerating}: (1) Accelerating the DP algorithms using hardware accelerators and (2) Developing pre-alignment filtering heuristics that reduce the need for the DP algorithms, given an edit distance threshold.

Hardware accelerators include building aligners that use 1) multi-core and SIMD (single instruction multiple data) capable central processing units (CPUs), such as Parasail \citep{daily2016parasail}.
The classical DP algorithms can also be accelerated by calculating a bit representation of the DP matrix and processing its bit-vectors in parallel, such as Myers's bit-vector algorithm \citep{myers1999fast}. To our knowledge, Edlib \citep{vsovsic2017edlib} is currently the best-performing implementation of Myers's bit-vector algorithm. 
Other hardware accelerators include 2) graphics processing units (GPUs), such as GSWABE \citep{liu2015gswabe}, 3) field-programmable gate arrays (FPGAs), such as FPGASW \citep{fei2018fpgasw}, or 4) {processing-in-memory architectures that enable performing computations inside the memory chip and alleviate the need for transferring the data to the CPU cores, such as} GenASM~\citep{senolcalimicro2020}. However, many of these efforts either simplify the scoring function as in Edlib, or only take into account accelerating the computation of the DP matrix without performing the backtracking step as in \citep{ chen2014accelerating}. Different and more sophisticated scoring functions are typically needed to better quantify the similarity between two sequences \citep{wang2011comparison}. The backtracking step involves unpredictable and irregular memory access patterns, which pose a difficult challenge for efficient hardware implementation.

Pre-alignment filtering heuristics aim to quickly eliminate some of the dissimilar sequences before using the computationally-expensive optimal alignment algorithms.
Existing pre-alignment filtering techniques are either: 1) slow and they suffer from a limited sequence length ($\leq128 bp$), such as SHD \citep{xin2015shifted}, or 2) inaccurate after some edit distance threshold, such as GateKeeper \citep{alser2017gatekeeper} and MAGNET \citep{alser2017magnet1}.
{Highly-parallel filtering can also be achieved using processing-in-memory architectures, as in} GRIM-Filter~\citep{kim2018grim}.
Shouji \citep{alser2019shouji} is currently the best-performing FPGA pre-alignment filter in terms of both accuracy and execution time.
%We provide full descriptions of these two key directions of accelerating sequence alignment in \textcolor{red}{Supplementary Material, Section 5}.

%\textbf{Hardware accelerators are becoming increasingly popular for speeding up the computationally expensive sequence aligners} \citep{aluru2014review, sandes2016parallel, al2017survey, ng2017reconfigurable}. 
%Hardware accelerators include multi-core and SIMD (single instruction multiple data) capable central processing units (CPUs), graphics processing units (GPUs), and field-programmable gate arrays (FPGAs). 
%The classical DP algorithms are typically accelerated by computing only the necessary regions (i.e., diagonal vectors) of the DP matrix rather than the entire matrix, as proposed in Ukkonen's banded algorithm \citep{ukkonen1985algorithms}. 
%The number of the diagonal bands required for computing the DP matrix is 2$E$+1, where $E$ is a user-defined edit distance threshold. The banded algorithm is still beneficial even with its recent sequential implementations as in Edlib \citep{vsovsic2017edlib}. 
%\alser{Edlib is implemented in C for standard CPUs and it calculates the banded Levenshtein distance}. 
%Parasail \citep{daily2016parasail} exploits both Ukkonen's banded algorithm and SIMD-capable CPUs to compute a banded alignment for a sequence pair with a user-defined scoring function. SIMD instructions offer significant parallelism to the matrix computation by executing the same vector operation on multiple operands at once. Multi-core architecture of CPUs and GPUs provides the ability to compute alignments of many sequence pairs independently and concurrently \citep{georganas2015meraligner, liu2015gswabe}. GSWABE \citep{liu2015gswabe} exploits GPUs (Tesla K40) for a highly-parallel computation of global alignment with a user-defined scoring function. CUDASW++ 3.0 \citep{liu2013cudasw++} exploits the SIMD capability of both CPUs and GPUs (GTX690) to accelerate the computation of the Smith-Waterman algorithm with a user-defined scoring function. CUDASW++ 3.0 provides only the optimal score, not the optimal alignment (i.e., no backtracking step). Other designs, for instance FPGASW \citep{fei2018fpgasw}, exploit the very large number of hardware execution units in FPGAs (Xilinx VC707) to form a linear systolic array \citep{kung1982systolic}. Each execution unit in the systolic array is responsible for computing the value of a single entry of the DP matrix. The systolic array computes a single vector of the matrix at a time. The data dependencies between the entries restrict the systolic array to computing the vectors sequentially (e.g., top-to-bottom, left-to-right, or in an anti-diagonal manner). FPGA accelerators seem to yield the highest performance gain compared to the other hardware accelerators \citep{fei2018fpgasw, waidyasooriya2016hardware, banerjee2019asap, chen2016spark}. However, many of these efforts either simplify the scoring function, or only take into account accelerating the computation of the DP matrix without performing the backtracking step as in \citep{liu2013cudasw++, nishimura2017accelerating, chen2014accelerating}. Different and more sophisticated scoring functions are typically needed to better quantify the similarity between two sequences \citep{wang2011comparison}. The backtracking step involves unpredictable and irregular memory access patterns, which poses a difficult challenge for efficient hardware implementation.

%\textbf{Pre-alignment filtering heuristics aim to quickly eliminate some of the dissimilar sequences before using the computationally expensive optimal alignment algorithms}. There are a few existing filtering techniques such as the Adjacency Filter \citep{xin2013accelerating}, which is implemented for standard CPUs as part of FastHASH \citep{xin2013accelerating}. SHD \citep{xin2015shifted} is a SIMD-friendly bit-vector filter that provides higher filtering accuracy compared to the Adjacency Filter. GRIM-Filter \citep{kim2018grim} exploits the high memory bandwidth and the logic layer of 3D-stacked memory to perform highly-parallel filtering in the DRAM chip itself. GateKeeper \citep{alser2017gatekeeper}, MAGNET \citep{alser2017magnet1}, and Shouji \citep{alser2019shouji} are designed to utilize the large amounts of parallelism offered by FPGA architectures. Shouji \citep{alser2019shouji} is currently the best-performing pre-alignment filter in terms of accuracy and execution time.

Our \textbf{goal} in this work is to significantly reduce the time spent on calculating the sequence alignment of \emph{both short and long sequences} using very fast and accurate pre-alignment filtering. 
To this end, we introduce \emph{SneakySnake}, a highly parallel and highly accurate pre-alignment filter that works on {\emph{modern} high-performance computing architectures such as CPUs, GPUs, and FPGAs}. \alser{The \textbf{key idea} of SneakySnake is to provide {a} highly-accurate pre-alignment filtering algorithm by reducing the ASM problem to the \emph{single net routing} (SNR) problem \citep{lee1976use}. 
The SNR problem is to find the shortest routing path that interconnects two terminals on the boundaries of VLSI chip layout {while passing} through the minimum number of obstacles.
Solving the SNR problem is faster than solving the ASM problem, as calculating the routing path after facing an obstacle is independent of the calculated path before this obstacle.}
This provides two key benefits. 1) It obviates the need for using computationally costly DP algorithms to keep track of the subpath that provides {the} optimal {solution} (i.e., {the one} with the least possible routing cost).
2) The independence {of} the subpaths allows for solving many SNR subproblems in parallel by judiciously leveraging the parallelism-friendly architecture of modern FPGAs and GPUs to greatly speed up the SneakySnake algorithm.

The \textbf{contributions} of this paper are as follows: 
\begin{itemize}
\item We introduce SneakySnake, 
%a highly parallel and highly accurate pre-alignment filter, which \alser{1) reduces the ASM problem to the SNR problem, 2) efficiently and optimally solves the SNR problem, 3) uses the SNR solution to quickly identify dissimilar sequences, and 4) avoids performing computationally costly sequence alignment for such dissimilar sequences.}
the fastest and most accurate pre-alignment filtering mechanism to date that greatly enables the speeding up of genome sequence alignment while preserving its accuracy. We demonstrate that the SneakySnake algorithm is 1) correct and optimal in solving the SNR problem and 2) it runs in linear time with respect to sequence length and edit distance threshold.
\item We demonstrate that the SneakySnake algorithm significantly improves the accuracy of pre-alignment filtering by up to four orders of magnitude compared to Shouji, GateKeeper, and SHD.
\item We provide, to our knowledge, the \emph{first universal} pre-alignment filter for {CPUs, GPUs, and FPGAs}, by having software as well as software/hardware co-designed versions {of SneakySnake}.
\item We demonstrate, using short sequences, that SneakySnake accelerates Edlib and Parasail by up to 37.7$\times$ and 43.9$\times$ (>12$\times$ on average), respectively, {with} its CPU implementation, and by up to 413$\times$ and 689$\times$ (>400$\times$ on average), respectively, {with} {FPGA and GPU} acceleration. 
{We also demonstrate, using long sequences, that SneakySnake accelerates Parasail by up to 979$\times$ (276.9$\times$ on average).}
\item We demonstrate that {the CPU implementation of} SneakySnake accelerates the sequence alignment of minimap2 \citep{li2018minimap2}, a state-of-the-art read mapper, {by up to 6.83$\times$ and  91.7$\times$ using short and long sequences, respectively}.
%Unlike most existing works that aim to accelerate sequence alignment, SneakySnake does not sacrifice \emph{any} of the aligner capabilities (i.e., scoring and backtracking), as \alser{it does} not replace the alignment step. 
%\item We demonstrate that both SneakySnake and Snake-on-Chip reduce the execution time of state-of-the-art read mapper, mrFAST \citep{alkan2009personalized}, by up to 2.93$\times$ and 7.57$\times$, respectively.
\end{itemize}
%\vspace{-28 pt}